\documentclass[11pt]{amsart}
\usepackage[a4paper,margin=2.6cm]{geometry}
\usepackage{amssymb,amsmath,booktabs,longtable,enumitem}
\usepackage[hidelinks]{hyperref}

\newtheorem{theorem}{Theorem}[section]
\newtheorem{proposition}[theorem]{Proposition}
\newtheorem{lemma}[theorem]{Lemma}
\newtheorem{corollary}[theorem]{Corollary}
\newtheorem{conjecture}[theorem]{Conjecture}
\theoremstyle{definition}
\newtheorem{definition}[theorem]{Definition}
\theoremstyle{remark}
\newtheorem{remark}[theorem]{Remark}
\newtheorem{observation}[theorem]{Observation}

\newcommand{\Q}{\mathbb{Q}}
\newcommand{\R}{\mathbb{R}}
\newcommand{\Z}{\mathbb{Z}}
\newcommand{\Hb}{\mathbb{H}}
\newcommand{\OK}{\mathcal{O}_K}
\newcommand{\GL}{\mathrm{GL}}
\newcommand{\PGL}{\mathrm{PGL}}
\newcommand{\PSL}{\mathrm{PSL}}
\newcommand{\SL}{\mathrm{SL}}
\newcommand{\Reg}{\mathrm{Reg}_K}
\newcommand{\Aut}{\mathrm{Aut}}
\newcommand{\gK}{\gamma_{K,2}}
\newcommand{\nperf}{n_{\mathrm{perf}}}
\newcommand{\nex}{n_{\mathrm{ext}}}
\newcommand{\PROVED}{\textsc{[proved]}}
\newcommand{\CERT}{\textsc{[certified]}}
\newcommand{\CONJ}{\textsc{[conjecture]}}
\newcommand{\REFQ}{\textbf{[REF?]}}

\title[Perfect binary Humbert forms, multiplicative normalisation]{Perfect and extreme binary Humbert forms over real quadratic fields
in the multiplicative normalisation:\\ an enumeration for 157 fields, patterns, and next steps}
\date{Working notes, November 2026}

\begin{document}

\begin{abstract}
For every real quadratic field $K=\Q(\sqrt D)$, $D\le 300$ squarefree, whose Hermite--Humbert
constant $\gK$ (Icaza's multiplicative normalisation) has been certified in this project ---
157 of the 182 fields --- we list all classes of perfect binary Humbert forms modulo
$\GL_2(\OK)$ and scaling, and decide eutaxy for each.  There are $55444$ perfect classes, of
which $15313$ are extreme, $80$ weakly eutactic and $40051$ not eutactic; in every field the
largest $\gamma$ of a perfect form equals the certified $\gK$.  We describe the enumeration
(a branch-and-bound over the fundamental region with certified discarding rules, followed by
a local Newton solve), prove that it is complete under one explicitly stated numerical
hypothesis, and discuss how to remove that hypothesis.  The data suggest three regularities:
the number of perfect classes is proportional to $\zeta_K(-1)$, i.e.\ to the volume of the
Hilbert modular surface, with no regulator factor (in contrast with the additive
normalisation); $\gK$ is described to within about $20\%$ (residual standard deviation $0.18$ in $\log$) by $d_K^{1.17}\Reg^{-0.61}$, and
we give an elementary proof, for $h_K=1$, of the case needed here of Thunder's lower bound
$\gK\ge 4\pi^2\zeta_K(-1)/(h_K\sqrt{d_K}\,\Reg)$, which the data satisfy even without the
factor $1/h_K$; and forms
with an automorphism of order $3$ have $\tfrac34\gamma$ equal to a norm from $\Z[\zeta_3]$,
reflecting a CM structure.  Every claim is labelled \PROVED, \CERT{} (certified or exhaustive
computation, range stated) or \CONJ.
\end{abstract}

\maketitle
\tableofcontents

\section{Introduction}\label{sec:intro}

Let $K=\Q(\sqrt D)$ be a real quadratic field with discriminant $d=d_K$, ring of integers
$\OK$, fundamental unit $\varepsilon_0>1$, regulator $\Reg=\log\varepsilon_0$ and class number
$h=h_K$.  A \emph{binary Humbert form} is a pair $S=(S_1,S_2)$ of positive definite real
$2\times2$ matrices, $S_j$ attached to the real embedding $x\mapsto x^{(j)}$.  Following
Icaza and \cite{BCIO,PW05,PW09} we use the \emph{multiplicative} normalisation
\begin{equation}\label{eq:mult}
S[v]=S_1[v^{(1)}]\,S_2[v^{(2)}],\quad m(S)=\min_{0\ne v\in\OK^2}S[v],\quad
\gamma(S)=\frac{m(S)}{\sqrt{\det S_1\det S_2}},\quad \gK=\sup_S\gamma(S).
\end{equation}
This is a different theory from the \emph{additive} normalisation
$S_1[v^{(1)}]+S_2[v^{(2)}]$ of Koecher, Leibak and Gunnells--Yasaki
\cite{Leibak,LeibakThesis,GY}: the minimum is a product rather than a sum, so it is not a
linear function of $S$, and there is no Ryshkov polyhedron.  The two are never mixed below;
the additive theory is the subject of the companion project \texttt{pav10/perfect\_forms},
which we only quote for comparison.

Known values.  Cohn \cite{Cohn1} computed the lowest points of the Hilbert modular
fundamental domain for $\Q(\sqrt2)$, $\Q(\sqrt3)$, $\Q(\sqrt5)$, which give $\gK$ there, and
these values were established rigorously in \cite{BCIO}.  Pohst and Wagner obtained
$\Q(\sqrt{13})$ and the extreme forms of $\Q(\sqrt{17})$ having a unimodular pair of minimal
vectors, without deciding whether these are all \cite{PW05}; then $\Q(\sqrt6)$ ($\gK=5$) and
$\Q(\sqrt{21})$ ($\gK=16/3$), and conjectured the values for $\Q(\sqrt7)$ and $\Q(\sqrt{11})$
\cite{PW09}.  Coulangeon, Icaza and O'Ryan computed $\gK$ for the complex cubic field
$\Q(\rho)$, $\rho^3-\rho^2+1=0$ \cite{CIO}, and note that every value known then was obtained by
constructing the possible sets of minimal vectors of perfect forms.  In this project it was certified for $157$ fields with $D\le300$
(branch-and-bound over the fundamental region, every certificate re-checked with outward
rounding; see the project README).  The present notes report a second computation on top of
the first: the complete list of \emph{perfect} forms (in Coulangeon's sense
\cite{Coulangeon}) and their \emph{eutaxy}, for all $157$ fields.  The extreme forms (perfect
and eutactic) are the local maxima of $\gamma$; the global maximum is $\gK$.

\subsection*{Summary}
\begin{enumerate}[leftmargin=*]
\item \emph{Algorithm} (\S\ref{sec:algo}).  Voronoi's neighbour algorithm is not available in
the multiplicative setting (\S\ref{subsec:whynot}).  We enumerate instead: a branch-and-bound
over the region of highest points discards boxes with three certified rules, and the
remaining small boxes are solved locally.  Theorem~\ref{thm:complete} proves completeness
under a single hypothesis (H) on the local Newton solve, and \S\ref{subsec:improve} explains how
interval Newton would remove it.  The cost is proportional to $\zeta_K(-1)$
(\S\ref{subsec:cost}).
\item \emph{Results} (\S\ref{sec:results}).  $55444$ perfect classes in $157$ fields
(Table~\ref{tab:counts}); validation against the literature, against the certified $\gK$, and
by an independent re-run at half the box size.
\item \emph{Patterns} (\S\ref{sec:patterns}).  A volume law $\nperf\approx 7.9\,\zeta_K(-1)$
(Conjecture~\ref{conj:volume}); an empirical law and a proved lower bound for $\gK$
(Proposition~\ref{thm:lower}, after Thunder); CM structure of the symmetric forms (Observation~\ref{obs:cm}).
\item \emph{Next steps} (\S\ref{sec:next}) and the literature we need (\S\ref{sec:lit}).
\end{enumerate}

\subsection*{Conventions}
Labels: \PROVED{} --- proved here or in a cited source; \CERT{} --- established by a
certified or exhaustive computation in the stated range (for the enumeration of perfect forms
this means: complete under hypothesis (H) of \S\ref{subsec:correct}); \CONJ{} --- conjecture
or heuristic.  Fits are least-squares fits of logarithms; $R^2$ is the coefficient of
determination.  All numbers are produced by \texttt{analysis/patterns.py} from
\texttt{results/perfect/} and \texttt{results/gamma\_K2.tsv}.

\section{Setting}\label{sec:setting}

\subsection{Forms as points of $\Hb^2$}
Up to the two independent scalings $(S_1,S_2)\mapsto(\lambda_1S_1,\lambda_2S_2)$, which leave
$\gamma$ unchanged, a form is a point $z=(z_1,z_2)\in\Hb^2$, $z_j=x_j+iy_j$, via
$S_j\propto\frac1{y_j}\begin{pmatrix}1&x_j\\x_j&|z_j|^2\end{pmatrix}$.  For $v=(a,b)\in\OK^2$,
\[
f_v(z)=\frac{|a^{(1)}+b^{(1)}z_1|^2\,|a^{(2)}+b^{(2)}z_2|^2}{y_1y_2},\qquad
G(z)=\min_{v\ne0}f_v(z)=\gamma(S).
\]
$G$ is invariant under $\GL_2(\OK)$ and under the Galois swap $(z_1,z_2)\mapsto(z_2,z_1)$, and
$f_{\varepsilon v}=f_v$ for units $\varepsilon$, so minimal vectors come in \emph{unit classes}.
The space of forms modulo scaling is thus $\Hb^2$ itself (real dimension $4$), and forms modulo
equivalence are points of the Hilbert modular surface $\PGL_2(\OK)\backslash\Hb^2$.
(In the additive normalisation only the common scaling is free, and the space has dimension
$5$; see \S\ref{subsec:volume}.)

\subsection{Highest points and the search region}
For $h=1$ every orbit has a \emph{highest} point (maximal $y_1y_2$), and there $e_1=(1,0)$ is
a minimal vector, so $G(z)=f_{e_1}(z)=1/(y_1y_2)$ \cite[Prop.~1.1]{BCIO}.  Using translations
by $\OK$, $z\mapsto\varepsilon_0^kz$, $x\mapsto-x$ (all in $\PGL_2(\OK)$) and the Galois swap,
every orbit has a highest point in
\[
\mathcal R=\{x=r_1+r_2\omega,\ r_1\in[-\tfrac12,\tfrac12],\ r_2\in[0,\tfrac12];\
y_2/y_1\in[1,\varepsilon_0];\ y_1y_2\ge 2/d\},
\]
$\omega=(1+\sqrt D)/2$ or $\sqrt D$; the last condition is Cohn's bound $\gK\le d/2$
\cite{Cohn1,Cohn2}.  For $h>1$ the project uses one \emph{chart} per pair of ideal classes
$\{C,C^{-1}\}$: in the chart of an integral ideal $\mathfrak c$ of minimal norm in $C$ the
lattice is $\mathfrak c^{-1}\oplus\mathfrak c$, with the same region up to the translation lattice
(project README, ``Steinitz charts'').  We write $t=\frac12\log(y_1y_2)$,
$s=\frac12\log(y_2/y_1)$, and use box coordinates $(x_1,x_2,s,t)$; within a chart put
$N_v=f_v/f_{e_1}$, so that $N_v(z)=|a^{(1)}+b^{(1)}z_1|^2|a^{(2)}+b^{(2)}z_2|^2$.

\subsection{Perfection and eutaxy}
For a form $S$ and a minimal vector $v$ let
\[
X_v=\bigl(S_2[v^{(2)}]\,v^{(1)}v^{(1)T},\ S_1[v^{(1)}]\,v^{(2)}v^{(2)T}\bigr),
\]
the gradient of $S\mapsto S[v]$ in the $6$-dimensional space $\mathrm{Sym}_2(\R)^2$; $X_{\varepsilon v}=X_v$
for units.  $S$ is \emph{perfect} if the $X_v$ ($v$ minimal) span a space of dimension $5$,
the maximum possible \cite{Coulangeon,BCIO}; \emph{eutactic} if
$(S_1^{-1},S_2^{-1})=\sum_v\rho_vX_v$ with all $\rho_v>0$; \emph{weakly eutactic} if this
holds with $\rho_v\ge0$ but not with $\rho_v>0$; and \emph{extreme} (a strict local maximum of
$\gamma$) iff perfect and eutactic \cite{Coulangeon}.

\section{The enumeration}\label{sec:algo}

\subsection{Why not Voronoi's algorithm}\label{subsec:whynot}
In the additive theory $S\mapsto S_1[v^{(1)}]+S_2[v^{(2)}]$ is linear, the forms with minimum
$\ge1$ form a locally finite convex polyhedron (Ryshkov), perfect forms are its vertices, and
Voronoi's algorithm walks from vertex to neighbouring vertex
\cite{Voronoi,Leibak,GY}.  In the multiplicative theory the constraint $S[v]\ge1$ is
not linear (it is $\log S_1[v^{(1)}]+\log S_2[v^{(2)}]\ge0$), the region $\{m\ge1\}$ is not
convex, and the local structure at a perfect point --- although described by Coulangeon's
theory --- does not give a finite combinatorial ``neighbour'' step that is guaranteed to
succeed.  We implemented such a walk (\texttt{voronoi.py}: Lipschitz-certified steps along the
edges of the arrangement $\{f_v=G\}$); it works on generic vertices but fails on degenerate
faces, where more than the expected number of $f_v$ meet.  We therefore enumerate.

\subsection{The algorithm}\label{subsec:algo}
Fix a chart.  Let $\mathcal V$ be a finite set of unit classes containing every class $v$ with
$f_v\le d/2$ somewhere on $\mathcal R$ (it exists and is computed exactly as in the
$\gK$-certificate).  For a box $B$ of $(x_1,x_2,s,t)$-space with $t\le0$:
\begin{description}[leftmargin=*]
\item[Rule 3] discard $B$ if it misses the translation parallelogram;
\item[Rule 1] discard $B$ if $\sup_BN_v<1$ for some $v\in\mathcal V$ (then $e_1$ is not
minimal anywhere in $B$, so $B\cap\mathcal R'=\emptyset$, $\mathcal R'$ the set of points of
$\mathcal R$ at which $e_1$ is minimal);
\item[Rule 5] discard $B$ if fewer than $5$ classes $v\in\mathcal V$ (including $e_1$, for which
$N_{e_1}\equiv1$) satisfy $\inf_BN_v\le1$.
\end{description}
The bounds $\sup_BN_v$, $\inf_BN_v$ are the product bounds of the $\gK$-certificate (exact
interval bounds for $|a+bz|^2$ on a box, evaluated in floating point with a relative safety
margin $10^{-9}$).  Boxes that survive are bisected along the coordinate of largest
\emph{hyperbolic} width ($\Delta x_j/y_j$, $\Delta s$, $\Delta t$) until that width is
$\le w_{\min}=0.05$.  In a surviving leaf $B$, with candidate set
$C_B=\{v\ne e_1:\inf_BN_v\le1\}$:
\begin{enumerate}[leftmargin=*]
\item for every $4$-subset $\{v_1,\dots,v_4\}\subset C_B$, run damped Newton on
$\log N_{v_i}(q)=0$, $i=1,\dots,4$, in $q=(x_1,x_2,\log y_1,\log y_2)$, from the centre of
$B$ and three random points of $B$; accept the first limit that lies in $B$ (with a margin
of $1\%$ of the box size);
\item polish an accepted point to $50$ digits; compute the full set $M$ of minimal classes
\emph{independently of $\mathcal V$} by the exact short-vector enumeration; reject the point
if $e_1$ is not minimal; compute the rank of $\{X_v\}_{v\in M}$ and keep the point if it is
$5$;
\item decide eutaxy by a linear programme (and, when $|M|=5$, by solving for the unique
$\rho$ at $50$ digits).
\end{enumerate}
Finally the perfect points from all charts are grouped into $\GL_2(\OK)$-classes by
\cite[Alg.~11]{PW05} in exact arithmetic, and the stabiliser of each class in
$\GL_2(\OK)/\OK^\times$ is computed.  Large fundamental units are handled by cutting the
$s$-range into slabs.

\subsection{Correctness}\label{subsec:correct}

\begin{lemma}[\PROVED]\label{lem:height}
Let $z\in\mathcal R'$ be perfect.  Then $z$ has a minimal class $v=(a,b)$ with $b\ne0$, and
$y_1y_2\le1$, i.e.\ $t\le0$; in the chart of $\mathfrak c$, $y_1y_2\le N(\mathfrak c)^{-1}\le1$.
For $h=1$ every minimal class other than that of $e_1$ has $b\ne0$.
\end{lemma}
\begin{proof}
For a minimal $v=(a,0)$ we have $S_j[v^{(j)}]=a^{(j)2}S_j[e_1]$ and $S[v]=S[e_1]=m$, so
$X_v=\bigl(\tfrac{m}{S_1[e_1]}E_{11},\tfrac{m}{S_2[e_1]}E_{11}\bigr)=X_{e_1}$ ($E_{11}=e_1e_1^T$).
So classes with $b=0$ add nothing to the rank, and a perfect point (rank $5$) has a minimal class
with $b\ne0$.  For it,
$f_v\ge (b^{(1)}y_1)^2(b^{(2)}y_2)^2/(y_1y_2)=N(b)^2y_1y_2$, while $f_v=f_{e_1}=1/(y_1y_2)$;
hence $(y_1y_2)^2\le N(b)^{-2}$.  For $h=1$, $|N(b)|\ge1$; in a chart $b\in\mathfrak c$, so
$|N(b)|\ge N(\mathfrak c)\ge1$.  Finally, for $h=1$ and $b=0$, $f_v=N(a)^2f_{e_1}$ forces
$|N(a)|=1$, so $a$ is a unit and $v$ is in the class of $e_1$.
\end{proof}

\begin{lemma}[\PROVED]\label{lem:rank}
At $z\in\mathcal R'$ with minimal classes $M\ni e_1$, the rank of $\{X_v\}_{v\in M}$ is $5$ if
and only if the differentials $d\log N_v$, $v\in M\setminus\{e_1\}$, span the $4$-dimensional
cotangent space of $\Hb^2$ at $z$.  In that case some $4$ of them are linearly independent,
and $z$ is an isolated, nondegenerate solution of the corresponding $4\times4$ system
$\log N_{v_i}=0$.
\end{lemma}
\begin{proof}
For minimal $v$, $\langle X_v,(S_1,0)\rangle=\langle X_v,(0,S_2)\rangle=S[v]=m$, so all $X_v$
lie in an affine subspace $A$ of codimension $2$ not containing $0$.  Hence
$\operatorname{rank}\{X_v\}=1+\operatorname{rank}\{X_v-X_{e_1}\}$, and the differences lie in
the annihilator of the two scaling directions, i.e.\ in the cotangent space of the quotient
$\Hb^2$, of dimension $4$.  There $X_v-X_{e_1}$ is $m$ times $d\log(S[v]/S[e_1])=d\log N_v$.
So rank $5$ $\iff$ the $d\log N_v$ span a space of dimension $4$, and the Jacobian of a suitable
$4$-subset is invertible at $z$.
\end{proof}

\begin{lemma}[\PROVED, up to floating point]\label{lem:rules}
Rules 1, 3 and 5 never discard a box containing a perfect point $z\in\mathcal R'$ with $t\le0$.
Every minimal class of such a $z$ lies in $\mathcal V$ and in the candidate set of every
leaf containing $z$.
\end{lemma}
\begin{proof}
At $z$, $N_v(z)\ge1$ for all $v$ ($e_1$ minimal), so $\sup_BN_v\ge1$: rule 1 does not apply;
$z$ lies in the parallelogram: rule 3 does not apply; $z$ has at least $5$ minimal classes
(perfection needs rank $5$), each with $N_v(z)=1$, hence $\inf_BN_v\le1$: rule 5 does not
apply.  A minimal class has $f_v(z)=G(z)\le\gK\le d/2$, so it is in $\mathcal V$.  The bounds
are exact interval bounds evaluated in floating point with a safety margin; ``up to floating
point'' refers to this evaluation, which is the same as in the $\gK$-certificate (there
re-checked with outward rounding).
\end{proof}

\noindent
\textbf{Hypothesis (H).}  \emph{For every perfect point $z$ and every leaf $B\ni z$, Newton for
at least one $4$-subset of minimal classes of $z$ with invertible Jacobian (Lemma~\ref{lem:rank})
converges to $z$ from one of the four starting points, and no other solution in $B$ of the
same subset is reached first.}

\begin{theorem}[completeness, \PROVED{} under (H)]\label{thm:complete}
Under (H) the enumeration returns a representative of every $\GL_2(\OK)$-class of perfect
binary Humbert forms, and only perfect forms.
\end{theorem}
\begin{proof}
Every class has a highest point in some chart's region with $e_1$ minimal and, by
Lemma~\ref{lem:height}, $t\le0$.  The leaves and discarded boxes tile the box cover of the
region, and by Lemma~\ref{lem:rules} $z$ lies in a leaf $B$ whose candidate set contains its
minimal classes.  By (H) some $4$-subset produces $z$; the polished point has the minimal set
and rank computed independently of $\mathcal V$, so it is kept, and only points with rank $5$
are kept.  The classification is exact.
\end{proof}

\begin{remark}[what is not rigorous]\label{rem:nonrig}
(i) Hypothesis (H): Newton has a nonsingular solution in the box but convergence from the
chosen starts is not proved; the solve is local (box width $\le0.05$ in hyperbolic terms) and
started from $4$ points, and a failure would need all of them to fail.  (ii) The rank $5$ and
the eutaxy decisions are made at $50$ digits with tolerances ($10^{-25}$ relative for singular
values, $10^{-9}$ for the eutaxy margin), so a weakly eutactic class is weakly eutactic to within
$10^{-9}$.  (iii) The class list is checked against the certified $\gK$ only through its maximum.
\end{remark}

\subsection{Cost}\label{subsec:cost}
Over the $157$ fields the search used $7.2\cdot10^8$ boxes, $5.3\cdot10^7$ leaves and
$1.5\cdot10^9$ Newton runs.  The number of boxes grows like $\zeta_K(-1)^{1.02}$ ($R^2=0.91$,
the $127$ fields with $d\ge100$), the same law as the number of perfect classes
(\S\ref{subsec:volume}): the work is proportional to the volume of the Hilbert modular
surface, i.e.\ roughly to $d^{3/2}$.  The single most expensive field took $3.2$ hours on one
core ($D=263$).  Each class is found $3.1$ times on average (different charts, translates and
slabs), and $7.2\%$ of the classes only once.

\subsection{Making it robust and faster}\label{subsec:improve}
\begin{enumerate}[leftmargin=*]
\item \emph{Certified completeness.}  Replace (H) by an interval Newton (Krawczyk) test per
$4$-subset and leaf: either it proves that the subset has no solution in $B$, or that it has
exactly one; leaves where neither holds are bisected.  At a perfect point the Jacobian of
some subset is invertible (Lemma~\ref{lem:rank}), so near it the test succeeds after finitely
many bisections; degenerate (non-perfect) crossings are rejected by the rank test.  Together
with an interval check of the rank and of the eutaxy LP this makes the lists certified.
\item \emph{Fewer subsets.}  Most of the time goes to $4$-subsets with no common zero in the
leaf.  Pairwise and triple tests (interval bounds on $\log N_v-\log N_w$) discard most of them
before Newton.
\item \emph{Fewer copies.}  Restricting the region by the stabilisers of the cusps (and the
chart symmetries) would remove most of the $3.1$ copies per class.
\item \emph{Parallelism inside a field.}  Boxes are independent; splitting the $s$- or
$t$-range over processes removes the long tail (now one process per field).
\item \emph{A compiled kernel} for the box bounds and the batched Newton.  (The batched Newton
already gave a factor $\approx3$ over the per-subset loop.)
\item \emph{A multiplicative neighbour walk} with exact handling of degenerate faces would make
the cost proportional to the output; a mass formula over the cells (as in the additive census
of \texttt{perfect\_forms}) would then give an independent completeness check.
\end{enumerate}

\subsection{Is the algorithm new?}\label{subsec:new}
The ingredients are known: highest points and the floor of the fundamental domain
\cite{Cohn1,BCIO,PW09}, perfection and eutaxy in this setting \cite{Coulangeon}, the
equivalence test \cite[Alg.~11]{PW05}, and the idea of solving $S[v]=1$ for a few candidate
minimal vectors.  The earlier algorithms differ from ours as follows.
\begin{itemize}[leftmargin=*]
\item \cite[Alg.~9]{PW05} and \cite{BCIO,CIO} normalise $e_1,e_2\in M(S)$ (a \emph{unimodular
pair} of minimal vectors), enumerate globally all $3$-sets of candidates from a finite set $M$
determined by norm bounds \cite[Lemmas~7, 8]{PW05}, solve the polynomial systems $S[u_i]=1$
(three unknowns $b_1,b_2,c$), and keep the solutions that are perfect and eutactic.  This
finds the \emph{extreme} forms with a unimodular pair; that every extreme form has one must be
proved field by field (\cite[Lemma~16]{PW05} for $\Q(\sqrt{13})$; \cite[Props.~2.9, 2.10]{CIO}
via Lenstra's constant), and for $\Q(\sqrt{17})$ it was left open.
\item \cite[\S2]{PW09} steps through Cohn's description of the fundamental domain with a
prescribed increment, collects the minimal vectors met in a set $M$, and solves the systems
for all $4$-sets of $M$; there is no criterion that $M$ contains all the vectors needed
\cite[p.~940]{PW09}.
\item Here only $e_1$ is normalised (it is minimal at the highest point, for any class
number, chart by chart), so no unimodular-pair hypothesis is needed; the $4$-subsets are
formed \emph{locally}, in boxes of hyperbolic width $\le0.05$ that survive certified
discarding rules, so that the candidate set provably contains the minimal classes of every
perfect point in the box (Lemma~\ref{lem:rules}); and all perfect forms are kept, not only the
eutactic ones.
\end{itemize}
To our knowledge this combination, and complete lists of multiplicative perfect forms for more
than a few fields, are new.  We have not seen Wagner's tables \cite{PW09} (cited there as a web
page) or later work by the same authors.

\section{Results}\label{sec:results}

\begin{table}[ht]
\caption{Totals over the $157$ fields with certified $\gK$, $D\le300$ \CERT.}\label{tab:counts}
\begin{tabular}{lr}
\toprule
perfect classes & $55444$\\
\quad extreme (perfect $+$ eutactic) & $15313$\\
\quad weakly eutactic (in $52$ fields) & $80$\\
\quad not eutactic & $40051$\\
fields with one perfect class & $D=2,3,5$\\
most perfect classes & $D=274$: $1467$\\
median fraction extreme & $0.332$\\
\bottomrule
\end{tabular}
\end{table}

\subsection{Validation}\label{subsec:validation}
\begin{itemize}[leftmargin=*]
\item \CERT{} In every one of the $157$ fields the largest $\gamma$ of a perfect class equals the
certified $\gK$ (relative difference $\le10^{-12}$).
\item \CERT{} Agreement with the literature.  $D=2,3,5$: one class each \cite{BCIO}
($|\Aut|=10$ for $D=5$).  $D=13$: two perfect classes, both extreme, $\gamma=(84+6\sqrt{91})/35
=\sqrt{(1476+144\sqrt{91})/175}$ (critical) and $\gamma=4.00045\ldots$ (five minimal classes),
the forms in rows $2$ and $1$ of the table of \cite[p.~916]{PW05}; the form of row $3$ has
perfection rank $<5$ there and is correctly absent.  $D=17$: exactly the two extreme forms of
\cite[pp.~918--919]{PW05}, $\gamma=\sqrt{(784+128\sqrt{34})/225}=2.6079893\ldots$ and
$(40+8\sqrt{85})/45=\sqrt{(1408+128\sqrt{85})/405}=2.5279190\ldots$, and no other perfect
class; so the question left open there is settled and $\gamma_{\Q(\sqrt{17}),2}=2.6079893\ldots$
(the $\gK$-certificate confirms the value).  $D=6$, $21$: $\gK=5$, $16/3$ as in \cite{PW09};
$D=7$, $11$: the conjectured values $\sqrt{(5584+640\sqrt{70})/605}=(100+16\sqrt{70})/55$ and
$28/3$ of \cite[p.~946]{PW09} are the certified ones.  $D=10$: $7$ perfect, $5$ extreme (test
suite).  Where the experimental
neighbour walk completed, it found the same classes.
\item \CERT{} Independent re-run: the fields with $D\le 77$ listed in the file
\texttt{robustness.txt} of \texttt{results/perfect/} were recomputed with half the leaf width
($w_{\min}=0.025$, hence different leaves and Newton starts); the class lists (value of
$\gamma$ to $25$ digits, number of minimal classes, eutaxy, $|\Aut|$) agree in every field.
\item An audit after the short-vector fix (project history): one record changed ($D=46$ had a
spurious duplicate class with one minimal vector missed; now $84$ perfect, $39$ extreme).
\end{itemize}

\subsection{The table}
Table~\ref{tab:all} in the appendix lists $\nperf$ and $\nex$ for every field; the file
\texttt{D$\langle D\rangle$.json} in \texttt{results/perfect/} contains every class (value of $\gamma$,
closed form when PSLQ finds one, minimal vectors, rank, eutactic coefficients, $|\Aut|$, a point).

\section{Patterns}\label{sec:patterns}

\subsection{A volume law for the number of perfect forms}\label{subsec:volume}
\begin{observation}[\CERT, $157$ fields]\label{obs:volume}
Over the $127$ fields with $d\ge100$ the log-log slope of $\nperf$ against $\zeta_K(-1)$ is
$1.017$ ($R^2=0.984$); against $\zeta_K(-1)\Reg$ it is $0.807$ ($R^2=0.892$), against $d$ it is
$1.525$ ($R^2=0.973$), against $h\Reg$ $1.604$ ($R^2=0.54$).  For the $98$ fields with $d>200$,
$\nperf/\zeta_K(-1)$ has median $7.90$, quartiles $7.37$, $8.56$ and range $[5.96,10.86]$; the
mass $\sum1/|\Aut|$ behaves the same way (median $7.74$).  For the extreme forms the slope
against $\zeta_K(-1)$ is $0.823$.
\end{observation}

\begin{conjecture}[\CONJ]\label{conj:volume}
$\nperf(K)\asymp\zeta_K(-1)$ as $d_K\to\infty$, with constants independent of $\Reg$ and $h_K$.
\end{conjecture}

Heuristic.  By Hirzebruch's formula \REFQ{} the orbifold Euler characteristic of
$\PSL_2(\OK)\backslash\Hb^2$ is $2\zeta_K(-1)$, and its hyperbolic volume (measure
$\prod_j dx_j\,dy_j/y_j^2$) is $8\pi^2\zeta_K(-1)$; $\zeta_K(-1)\asymp d^{3/2}L(2,\chi_d)$.  In the
multiplicative normalisation the space of forms modulo scaling is exactly $\Hb^2$, so a
locally finite configuration of perfect points with bounded local density has
$\asymp\mathrm{vol}$ classes.  In the additive normalisation the separate scalings of $S_1$
and $S_2$ are not symmetries; the space has one more dimension, the log-unit direction, and
\texttt{perfect\_forms} (FINDINGS \S3, additive, $d\le100$) finds
$n\asymp\zeta_K(-1)\Reg$ there; for unary forms Yasaki's bound \cite{Yasaki} relates the count
to the period of a continued fraction, i.e.\ to $\Reg$.  The three laws are consistent with
``number of perfect forms $\asymp$ volume of the space of forms modulo $\GL_n(\OK)$ and the
symmetries of the normalisation''.

The constant $7.9$ is not explained.  An Euler-characteristic argument would need the full cell
structure (edges and higher cells), which the multiplicative theory does not provide
(\S\ref{subsec:whynot}); a neighbour walk with cells would turn the heuristic into a mass
formula and a completeness certificate.

\subsection{The Hermite--Humbert constant}\label{subsec:gamma}

\begin{observation}[\CERT, $157$ fields]\label{obs:gamma}
A least-squares fit gives
\[
\log\gK\approx 1.172\log d-0.614\log\Reg-1.711\qquad(R^2=0.977,\ \text{residual s.d. } 0.18),
\]
against $R^2=0.923$ for $\log d$ alone; adding $\log h$ changes nothing (coefficient $-0.05$).
$\gK/d$ ranges from $0.068$ ($D=46$) to $0.358$ ($D=5$), well below Cohn's $\tfrac12$.  The
low values of $\gK/d$ for $D\equiv1\pmod8$ ($2$ split) are explained by the regulator: the
residuals of the fit have mean $-0.04$ ($2$ split), $-0.03$ ($2$ inert), $+0.02$ ($2$
ramified).
\end{observation}

Known bounds.  Upper: $\gK<d/2$ (Cohn, as quoted in \cite{PW05,PW09}; see
\cite{Cohn1,Cohn2}) and $\gK\le16d/\pi^2$ (\cite[Thm.~1]{Icaza}, $n=2$, $r=2$, $s=0$).  Lower:
$\gK\ge B(K)$, the least norm of a nontrivial integral ideal, whenever Lenstra's constant
satisfies $L(K)<4$ \cite[Prop.~4.1]{CIO}; for quadratic $K$, $B(K)\le4$.  A lower bound
linear in $\sqrt{d_K}$ is attributed to a preprint of Baeza in \cite[p.~21]{Icaza}
(not seen).  The strongest is the Minkowski--Hlawka-type bound of Thunder
\cite[Thm.~2]{Thunder}: with $n=2$, $d=1$ (his notation), $r=2$, $s=0$, $w=2$ it reads
$\gamma_2(K)^{2}\ge\zeta_K(2)\,d_K/(\pi^2h_K\Reg)$, and $\gamma_2(K)^{[K:\Q]}\le\gK$ (his height
is absolute, $H_A=\|A\cdot\|_{\mathbb A}^{1/[K:\Q]}$, and for $h_K>1$ his minimum divides by the
norm of the content ideal, so it can only be smaller).  With $\zeta_K(-1)=d^{3/2}\zeta_K(2)/(4\pi^4)$:
\begin{equation}\label{eq:thunder}
\gK\ \ge\ \frac{\zeta_K(2)\,d_K}{\pi^2h_K\Reg}\ =\ \frac{4\pi^2\zeta_K(-1)}{h_K\sqrt{d_K}\,\Reg}.
\end{equation}
For $h_K=1$ the following elementary argument gives the same bound; it explains why the
regulator enters: the cusp has length $\Reg$, so a fixed volume spread over a longer cusp
forces the floor of the fundamental domain (height $1/\gK$) up.

\begin{proposition}[\PROVED{} for $h_K=1$, given Hirzebruch's volume formula; the case $h_K=1$ of
\cite{Thunder}]\label{thm:lower}
$\gK\ge 4\pi^2\zeta_K(-1)/(\sqrt{d_K}\,\Reg)$.
\end{proposition}
\begin{proof}
Every orbit has a highest point $z\in\mathcal R$, and there $G(z)=1/(y_1y_2)\le\gK$, so the set
$\mathcal R(Y)=\{z\in\mathcal R:\ y_1y_2\ge Y\}$ with $Y=1/\gK$ meets every orbit of the group
$\Gamma$ generated by $\PGL_2(\OK)$ and the Galois swap.  Hence its hyperbolic volume is at least
that of $\Gamma\backslash\Hb^2$.  The $x$-part of $\mathcal R$ has area $\sqrt d/2$ (half a period
parallelogram of $\OK$, of area $\sqrt d$).  With $p=\log(y_1y_2)$, $q=\log(y_2/y_1)$,
$dy_1dy_2/(y_1y_2)^2=\tfrac12e^{-p}\,dp\,dq$ and $q\in[0,\Reg]$, so
$\mathrm{vol}\,\mathcal R(Y)=\frac{\sqrt d}2\cdot\frac{\Reg}{2}\cdot\frac1Y=\frac{\sqrt d\,\Reg\,\gK}4$.
On the other hand $[\PGL_2(\OK):\PSL_2(\OK)]=|\OK^\times/\OK^{\times2}|=4$ and the swap has index
$2$, so $\mathrm{vol}(\Gamma\backslash\Hb^2)=8\pi^2\zeta_K(-1)/8=\pi^2\zeta_K(-1)$ (Hirzebruch's
formula: the Euler characteristic of $\PSL_2(\OK)\backslash\Hb^2$ is $2\zeta_K(-1)$, and the
volume is $4\pi^2$ times it \REFQ).
\end{proof}

\begin{observation}[\CERT, $157$ fields]
All these bounds hold in the data.  $\gK$ divided by \eqref{eq:thunder} lies in $[1.46,14.5]$
(minimum at $D=5$, then $D=2$, $17$); for $h_K>1$ the ratio is $\ge4.35$, and even without the
factor $1/h_K$ it is $\ge2.17$.  $\gK/B(K)\ge1.053$ for $D\ne5$ (nearly attained at $D=2$).
$\gK/(d/2)\le0.72$ and $\gK/(16d/\pi^2)\le0.22$.  The ratio to \eqref{eq:thunder} grows with $d$
(log-log slope of $\gK$ against $\zeta_K(-1)/(\sqrt d\Reg)$: $1.27$).
\end{observation}

\begin{remark}
Since $\zeta_K(2)\in[\zeta(4),\zeta(2)^2]$, \eqref{eq:thunder} gives
$\gK\gg d/(h\Reg)$, and $h\Reg\ll\sqrt d\log d$ by the class number formula, so
$\gK\gg\sqrt d/\log d$ for all real quadratic $K$; with $\gK<d/2$ this brackets the growth.  The
data sit near $d^{1.17}\Reg^{-0.61}$ and do not depend on $h$.  Questions: is $\gK\asymp d/\Reg$
(the fit says no)?  Can the factor $1/h_K$ be removed from \eqref{eq:thunder} (the data say
yes, with room to spare)?  Is $\limsup\gK/d<\tfrac12$?
\end{remark}

\subsection{Minimal vectors and automorphisms}\label{subsec:aut}
\begin{observation}[\CERT]
Of the $55444$ perfect classes, $44643$ ($80.5\%$) have exactly $5$ minimal classes, the minimum
for perfection (Coulangeon's bound $\#M(S)\ge rn(n+1)/2+sn^2-(r+s-1)=5$, \cite[(2--1)]{CIO},
\cite[Cor.~4]{PW05}); $8472$ have $6$; the maximum is $24$ ($D=263$).  The stabiliser in
$\GL_2(\OK)/\OK^\times$ is trivial for $53228$ classes, of order $2$ for $2106$, and of order
$3,4,6,10,24$ for $4,64,40,1,1$ classes.
\end{observation}

\begin{observation}[\CERT, $157$ fields]\label{obs:cm}
All $45$ classes whose stabiliser has order divisible by $3$ have $\tfrac34\gamma$ equal to an
integer of the form $a^2+ab+b^2$ (a norm from $\Z[\zeta_3]$).  Of the $64$ classes with stabiliser of
order $4$, $51$ have $\gamma$ an integer of the form $a^2+b^2$.  The stabiliser of order $10$ is
that of the unique class of $\Q(\sqrt5)$ ($\gamma=4/\sqrt5$), and the one of order $24$ that of
$\Q(\sqrt3)$ ($\gamma=4$).
\end{observation}

\begin{proposition}[\PROVED]\label{prop:cm}
Let $g\in\GL_2(\OK)$ stabilise a form and have order $3$ in $\PGL_2(K)$.  Then
$g^2-tg+\delta=0$ with $t\in\OK$, $\delta=\det g=t^2$, and $\OK^2$ is a module over the order
$\OK[g]\cong\OK[\tfrac12 t(1+\sqrt{-3})]$ of the CM field $L=K(\sqrt{-3})$.
\end{proposition}
\begin{proof}
$g^3$ is scalar and $g$ is not, so the eigenvalues are $\lambda,\lambda\zeta_3$; then
$t=\lambda(1+\zeta_3)=-\lambda\zeta_3^2$ and $\delta=\lambda^2\zeta_3$, so $t^2=\delta$ and the
discriminant is $t^2-4\delta=-3t^2$; $t\ne0$ since $\delta$ is a unit, so $g=\tfrac t2(1+\sqrt{-3})$ in
$L=K(\sqrt{-3})$.
\end{proof}

Heuristically, a form fixed by $g$ is, in each embedding, invariant under a rotation of order $3$, hence a
multiple of the hexagonal form, and $S[v]$ becomes a norm form of $L$ on a rank-one $\OK[g]$-module.
Since $L=K(\sqrt{-3})\supset\Q(\sqrt{-3})$, norms from $L$ to $\Q$ factor through
$\Q(\zeta_3)$; we expect this to explain the shape $\frac43(a^2+ab+b^2)$ in general \CONJ{}
(note $\frac43=\gamma_2^2$, the square of the Hermite constant of the hexagonal lattice).  The same mechanism with $\zeta_4$, $\zeta_5$
($\sqrt5\in K$) and $\zeta_{12}$ ($\sqrt3\in K$) accounts for the other symmetric forms.  This is
the correspondence with quadratic (CM) extensions of $K$ asked about: forms with extra
automorphisms carry the structure of hermitian forms over CM extensions $L/K$
(Proposition~\ref{prop:cm} for order $3$), and we expect their $\gamma$ to be governed by norms
from $L$ \CONJ.

\subsection{Eutaxy}
The median fraction of extreme classes per field is $0.33$, against $0.97$ for perfect unary
forms over real quadratic fields (\texttt{perfect\_forms}, FINDINGS \S5); weakly eutactic classes
occur ($80$, in $52$ fields), unlike the unary case.  The number of extreme classes grows more
slowly than $\nperf$ (slope $0.82$ against $\zeta_K(-1)$).

\section{Next steps}\label{sec:next}

\subsection{Completing $D\le300$ and certifying the lists}
The $25$ fields with $D\le300$ whose $\gK$ is not certified (large units) are the obvious
gap; the enumeration itself would run on them (it does not need $\gK$, only Cohn's bound), but
its cost grows with the length of the cusp.  Certified completeness needs the interval-Newton
version of \S\ref{subsec:improve} (estimated cost: a small constant factor, since most leaves
have no solution and are rejected after one interval evaluation).

\subsection{Ternary forms over real quadratic fields}
The space of ternary forms modulo the two scalings is $(\SL_3(\R)/\mathrm{SO}_3)^2$, of
dimension $10$.  A box enumeration in dimension $10$ is not feasible: the number of leaves
scales like $w_{\min}^{-10}$ times the volume of $\PGL_3(\OK)\backslash(\SL_3(\R)/\mathrm{SO}_3)^2$,
which already for $\Q(\sqrt5)$ is far beyond the binary case.  Feasible routes: (a) a
multiplicative Voronoi walk (\S\ref{subsec:improve}, item 6), whose cost is proportional to
the output; (b) a search for $\gamma_{K,3}$ alone, by local maximisation from many seeds
together with structured candidates (the lattices whose trace form is similar to $E_6$, or
hermitian structures over CM extensions as in Proposition~\ref{prop:cm}), with a
branch-and-bound certificate only if a good candidate is found; (c) in the additive
normalisation the Voronoi algorithm works in any rank, and \texttt{perfect\_forms} is the place
for it.  We would start with (b) for $\Q(\sqrt2)$, $\Q(\sqrt3)$, $\Q(\sqrt5)$.

\subsection{Binary forms over totally real fields of higher degree}
For a totally real cubic field the space of binary forms modulo the three scalings is $\Hb^3$
(dimension $6$), $\gamma_{K,2}$ is again the reciprocal height of the floor, and every step of
\S\ref{sec:algo} generalises (products of three factors, $\varepsilon$ in a rank-$2$ unit
lattice, $6$-subsets for Newton).  The cost of the $\gK$-certificate and of the enumeration is
proportional to the volume ($\asymp\zeta_K(-1)$ again) times $w_{\min}^{-6}$ instead of $w_{\min}^{-4}$:
feasible for the smallest cubic fields ($d=49,81,148,169,\dots$), out of reach for large ones
without the improvements above.  Fields with complex places are also natural first targets: a
complex place contributes a factor $\Hb^3$, and \cite{CIO} gives a value to test against
($\Q(\rho)$, $\rho^3-\rho^2+1=0$, $d_K=-23$: $\gK\approx2.4684958820$) and two further fields
($x^3+x^2-x+1$, $d_K=-44$; $x^4+x^3+x^2+x+1$, $d_K=125$) where the same method applies.  The code is written for two embeddings throughout; a degree-$n$
version is a rewrite of \texttt{humbert.py}, \texttt{certify.py} and \texttt{perfect.py}, not a
patch.

\subsection{Further questions suggested by the data}
\begin{enumerate}[leftmargin=*]
\item Prove $\nperf\asymp\zeta_K(-1)$ (Conjecture~\ref{conj:volume}), perhaps first a lower bound by
a covering argument as in Proposition~\ref{thm:lower}.
\item Explain $\gK\approx 0.18\,d^{1.17}\Reg^{-0.61}$; remove the factor $1/h_K$ from \eqref{eq:thunder}.
\item Classify the symmetric forms by CM extensions $L/K$ (which $L$, which orders, which
ideal classes) and compute their $\gamma$ from norms in $L$.
\item Compare the multiplicative and additive extreme forms field by field: \texttt{perfect\_forms}
finds $\gamma_{\mathrm{add}}^4/(16\gamma_{\mathrm{mult}}^2)\in[1.56,46.1]$, never $1$.
\end{enumerate}

\section{Literature}\label{sec:lit}
\emph{Read for this version} (provided by the project owner): \cite{PW05,PW09,CIO,Icaza,Thunder}.
They fix the novelty statement (\S\ref{subsec:new}), the validation (\S\ref{subsec:validation}),
the known bounds and the attribution of the lower bound to Thunder (\S\ref{subsec:gamma}).
\emph{Still wanted} (cited above only as reported by others, or marked \REFQ):
\begin{enumerate}[leftmargin=*]
\item R.~Baeza, \emph{The volume of the space of Humbert reduced forms}, preprint (cited in
\cite{Icaza} for a lower bound linear in $\sqrt{d_K}$) --- the closest relative of
Proposition~\ref{thm:lower}.
\item S.~Ohno and T.~Watanabe \cite{OW} (bounds for $\gamma_{K,n}$, used in \cite{CIO}).
\item R.~Baeza and M.~I.~Icaza \cite{BI97,BI04} (Humbert--Minkowski constant; unimodularity of
minimal vectors).
\item M.~Wagner's thesis and the table cited in \cite{PW09}; any later enumeration of
multiplicative perfect forms.
\item G.~Claus \cite{Claus} (lowest points of three Hilbert modular groups); H.~E.~Ong \cite{Ong} and
A.~Ash \cite{Ash} (a different notion of perfection over real quadratic fields, \cite[p.~937]{PW09}).
\item F.~Hirzebruch and G.~van der Geer on Hilbert modular surfaces (volume, Euler characteristic,
cusps) and Siegel's formula for $\zeta_K(-1)$ \REFQ.
\item Perfect and extreme hermitian forms over CM fields (for Proposition~\ref{prop:cm}), e.g.\ the
generalisations of Voronoi's theorem \cite{OY}.
\end{enumerate}
The additive theory \cite{Leibak,LeibakThesis,GY} is in the companion project.

\appendix
\section{Counts per field}\label{app:table}
\begin{small}
\begin{longtable}{rrrrr|rrrrr|rrrrr}
\caption{Per field: $D$, $h_K$, $\gK$ (rounded), number of perfect and of extreme classes \CERT.}\label{tab:all}\\
\toprule
$D$ & $h$ & $\gK$ & perf & ext & $D$ & $h$ & $\gK$ & perf & ext & $D$ & $h$ & $\gK$ & perf & ext\\
\midrule
\endhead
2 & 1 & 2.106 & 1 & 1 & 87 & 2 & 76 & 210 & 52 & 187 & 2 & 105.4 & 810 & 207\\
3 & 1 & 4 & 1 & 1 & 89 & 1 & 10.45 & 28 & 10 & 190 & 2 & 100.8 & 619 & 171\\
5 & 1 & 1.789 & 1 & 1 & 91 & 2 & 33.98 & 210 & 75 & 194 & 2 & 146.1 & 721 & 228\\
6 & 1 & 5 & 3 & 3 & 93 & 1 & 20.27 & 17 & 7 & 195 & 4 & 190.9 & 658 & 139\\
7 & 1 & 4.252 & 4 & 4 & 95 & 2 & 101.3 & 241 & 72 & 197 & 1 & 38.56 & 67 & 31\\
10 & 2 & 8.706 & 7 & 5 & 97 & 1 & 9.204 & 31 & 19 & 202 & 2 & 119.1 & 832 & 164\\
11 & 1 & 9.333 & 8 & 7 & 101 & 1 & 20.43 & 31 & 12 & 203 & 2 & 157 & 671 & 170\\
13 & 1 & 4.035 & 2 & 2 & 102 & 2 & 65.78 & 266 & 87 & 205 & 2 & 48.08 & 88 & 30\\
14 & 1 & 11.11 & 11 & 6 & 103 & 1 & 43.52 & 378 & 119 & 206 & 1 & 149.1 & 886 & 220\\
15 & 2 & 12.55 & 13 & 11 & 105 & 2 & 25.6 & 36 & 10 & 209 & 1 & 17.33 & 95 & 38\\
17 & 1 & 2.608 & 2 & 2 & 106 & 2 & 49.65 & 303 & 107 & 210 & 4 & 161.4 & 714 & 226\\
19 & 1 & 8.056 & 16 & 9 & 107 & 1 & 89.33 & 291 & 99 & 213 & 1 & 44.95 & 72 & 39\\
21 & 1 & 5.333 & 2 & 2 & 109 & 1 & 9.542 & 26 & 17 & 215 & 2 & 193.3 & 819 & 183\\
22 & 1 & 10.8 & 23 & 16 & 110 & 2 & 90.91 & 238 & 70 & 218 & 2 & 129.9 & 729 & 184\\
23 & 1 & 17.33 & 25 & 16 & 111 & 2 & 64.51 & 340 & 113 & 219 & 4 & 196 & 899 & 208\\
26 & 2 & 23.55 & 36 & 18 & 113 & 1 & 12.07 & 43 & 23 & 221 & 2 & 57.78 & 108 & 52\\
29 & 1 & 5.792 & 3 & 3 & 114 & 2 & 82.67 & 358 & 97 & 222 & 2 & 161.4 & 822 & 220\\
30 & 2 & 28 & 38 & 17 & 115 & 2 & 101.3 & 303 & 69 & 223 & 3 & 148.7 & 906 & 245\\
31 & 1 & 11.31 & 37 & 17 & 118 & 1 & 50.13 & 431 & 120 & 226 & 8 & 207.7 & 976 & 280\\
33 & 1 & 4.571 & 3 & 3 & 119 & 2 & 100.8 & 373 & 111 & 227 & 1 & 169.3 & 801 & 248\\
34 & 2 & 34.01 & 78 & 38 & 122 & 2 & 129.4 & 301 & 60 & 229 & 3 & 50.7 & 95 & 33\\
35 & 2 & 33.12 & 41 & 15 & 123 & 2 & 81.73 & 368 & 104 & 230 & 2 & 241.3 & 844 & 160\\
37 & 1 & 5.736 & 4 & 4 & 129 & 1 & 12.31 & 51 & 21 & 231 & 4 & 218.7 & 1032 & 263\\
38 & 1 & 30.16 & 34 & 19 & 130 & 4 & 117.6 & 439 & 114 & 233 & 1 & 21.55 & 143 & 54\\
39 & 2 & 25 & 54 & 31 & 131 & 1 & 91.24 & 465 & 114 & 235 & 6 & 281.3 & 1037 & 250\\
41 & 1 & 5.321 & 10 & 5 & 133 & 1 & 25.35 & 49 & 33 & 237 & 1 & 52.05 & 98 & 36\\
42 & 2 & 36.63 & 70 & 27 & 134 & 1 & 91.66 & 524 & 135 & 238 & 2 & 137.5 & 1357 & 280\\
43 & 1 & 21.13 & 87 & 27 & 137 & 1 & 13.29 & 47 & 22 & 246 & 2 & 90.05 & 1091 & 363\\
46 & 1 & 12.5 & 84 & 39 & 138 & 2 & 103.4 & 365 & 93 & 247 & 2 & 106.6 & 1181 & 240\\
47 & 1 & 33.41 & 73 & 38 & 141 & 1 & 28.09 & 63 & 27 & 253 & 1 & 29.8 & 129 & 43\\
51 & 2 & 36.6 & 91 & 44 & 142 & 3 & 106.2 & 448 & 137 & 254 & 3 & 184 & 1075 & 322\\
53 & 1 & 12.42 & 9 & 1 & 143 & 2 & 138.4 & 339 & 76 & 255 & 4 & 244.1 & 1058 & 257\\
55 & 2 & 40.11 & 118 & 48 & 145 & 4 & 35.14 & 102 & 36 & 257 & 3 & 52.21 & 145 & 47\\
57 & 1 & 9.333 & 10 & 5 & 146 & 2 & 129.3 & 520 & 165 & 258 & 2 & 204.3 & 1074 & 262\\
58 & 2 & 40.1 & 123 & 34 & 149 & 1 & 26.8 & 50 & 14 & 263 & 1 & 179.5 & 1144 & 303\\
59 & 1 & 44.1 & 110 & 45 & 154 & 2 & 78.99 & 534 & 153 & 265 & 2 & 34.11 & 246 & 78\\
61 & 1 & 6.456 & 8 & 7 & 155 & 2 & 125.8 & 531 & 118 & 266 & 2 & 146.7 & 1012 & 250\\
62 & 1 & 52.08 & 111 & 47 & 157 & 1 & 20.82 & 62 & 26 & 267 & 2 & 133.2 & 1124 & 320\\
65 & 2 & 17.07 & 21 & 7 & 158 & 1 & 109.1 & 578 & 142 & 269 & 1 & 49.51 & 126 & 43\\
66 & 2 & 52 & 142 & 47 & 159 & 2 & 111.2 & 565 & 164 & 273 & 2 & 33.33 & 267 & 113\\
67 & 1 & 29.27 & 177 & 65 & 161 & 1 & 12.63 & 79 & 32 & 274 & 4 & 209.6 & 1467 & 273\\
69 & 1 & 13.16 & 15 & 7 & 165 & 2 & 41.72 & 52 & 21 & 277 & 1 & 34.4 & 151 & 67\\
70 & 2 & 52.23 & 133 & 31 & 167 & 1 & 129.3 & 477 & 152 & 278 & 1 & 201 & 1095 & 296\\
71 & 1 & 50.08 & 210 & 67 & 170 & 4 & 158.8 & 566 & 170 & 282 & 2 & 183.6 & 1157 & 276\\
73 & 1 & 4.977 & 17 & 15 & 173 & 1 & 35.02 & 48 & 29 & 285 & 2 & 66.22 & 122 & 61\\
74 & 2 & 47.85 & 143 & 37 & 174 & 2 & 124 & 523 & 176 & 287 & 2 & 199.4 & 1134 & 358\\
77 & 1 & 17.63 & 12 & 8 & 177 & 1 & 14.3 & 88 & 40 & 290 & 4 & 279.4 & 1125 & 289\\
78 & 2 & 57.48 & 162 & 50 & 178 & 2 & 110 & 760 & 217 & 291 & 4 & 208.6 & 1427 & 398\\
79 & 3 & 52.14 & 203 & 74 & 181 & 1 & 17.92 & 68 & 35 & 293 & 1 & 70.92 & 109 & 39\\
82 & 4 & 75.75 & 228 & 57 & 182 & 2 & 176.7 & 539 & 111 & 298 & 2 & 165.3 & 1431 & 349\\
83 & 1 & 62.3 & 181 & 64 & 183 & 2 & 109 & 611 & 177 & 299 & 2 & 198.5 & 1223 & 304\\
85 & 2 & 22.1 & 47 & 14 & 185 & 2 & 36.13 & 106 & 16 &  & & & & \\
86 & 1 & 58.86 & 271 & 84 & 186 & 2 & 101 & 647 & 139 &  & & & & \\
\bottomrule
\end{longtable}
\end{small}


\begin{thebibliography}{99}

\bibitem{BCIO}
R.~Baeza, R.~Coulangeon, M.~I.~Icaza and M.~O'Ryan, \emph{Hermite's constant for quadratic number
fields}, Experiment. Math. \textbf{10} (2001), 543--551.

\bibitem{Ash}
A.~Ash, \emph{Deformation retracts with lowest possible dimension of arithmetic quotients of
self-adjoint homogeneous cones}, Math. Ann. \textbf{225} (1977), 69--76.

\bibitem{BI97}
R.~Baeza and M.~I.~Icaza, \emph{On Humbert--Minkowski's constant for a number field}, Proc. Amer.
Math. Soc. \textbf{125} (1997), 3195--3202.

\bibitem{BI04}
R.~Baeza and M.~I.~Icaza, \emph{On the unimodularity of minimal vectors of Humbert forms}, Arch.
Math. \textbf{83} (2004), 528--535.

\bibitem{Claus}
G.~Claus, \emph{Die Randmannigfaltigkeiten und die ``tiefsten'' Punkte des Fundamentalbereichs f\"ur
drei Hilbertsche Modulgruppen}, Math. Ann. \textbf{176} (1968), 225--256.

\bibitem{Cohn1}
H.~Cohn, \emph{A numerical survey of the floors of various Hilbert fundamental domains}, Math.
Comp. \textbf{19} (1965), 594--605.

\bibitem{Cohn2}
H.~Cohn, \emph{On the shape of the fundamental domain of the Hilbert modular group}, Proc. Sympos.
Pure Math. \textbf{8} (1965), 190--202.

\bibitem{CIO}
R.~Coulangeon, M.~I.~Icaza and M.~O'Ryan, \emph{Lenstra's constant and extreme forms in number
fields}, Experiment. Math. \textbf{16} (2007), 455--462.

\bibitem{Icaza}
M.~I.~Icaza, \emph{Hermite constant and extreme forms for algebraic number fields}, J. London Math.
Soc. (2) \textbf{55} (1997), 11--22.

\bibitem{OW}
S.~Ohno and T.~Watanabe, \emph{Estimates of Hermite constants for algebraic number fields},
Comment. Math. Univ. St. Paul \textbf{50} (2001), 53--63.

\bibitem{Ong}
H.~E.~Ong, \emph{Perfect quadratic forms over real-quadratic number fields}, Geom. Dedicata
\textbf{20} (1986), 51--77.

\bibitem{Thunder}
J.~L.~Thunder, \emph{Higher-dimensional analogs of Hermite's constant}, Michigan Math. J. \textbf{45}
(1998), 301--314.

\bibitem{Coulangeon}
R.~Coulangeon, \emph{Vorono\"i theory over algebraic number fields}, Monogr. Enseign. Math. \textbf{37}
(2001), 147--162.

\bibitem{GY}
P.~E.~Gunnells and D.~Yasaki, \emph{Perfect forms over totally real number fields}, arXiv:0908.3130.

\bibitem{Leibak}
A.~Leibak, \emph{On additive generalization of Vorono\"i's theory for algebraic number fields},
Proc. Estonian Acad. Sci. Phys. Math. \textbf{54} (2005), 195--211.

\bibitem{LeibakThesis}
A.~Leibak, \emph{On additive generalisation of Vorono\"i's theory of perfect forms over algebraic
number fields}, PhD thesis, Tallinn University of Technology, 2006.

\bibitem{OY}
K.~Okuda and S.~Yano, \emph{A generalization of Vorono\"i's theorem to algebraic lattices},
J. Th\'eor. Nombres Bordeaux \textbf{22} (2010), 727--740.

\bibitem{PW05}
M.~E.~Pohst and M.~Wagner, \emph{On the computation of Hermite--Humbert constants for real quadratic
number fields}, J. Th\'eor. Nombres Bordeaux \textbf{17} (2005), 905--920.

\bibitem{PW09}
M.~E.~Pohst and M.~Wagner, \emph{On the computation of Hermite--Humbert constants: the algorithm
of Cohn revisited}, J. Algebra \textbf{322} (2009), 936--947.

\bibitem{Voronoi}
G.~Vorono\"i, \emph{Nouvelles applications des param\`etres continus \`a la th\'eorie des formes
quadratiques. Premier m\'emoire}, J. Reine Angew. Math. \textbf{133} (1908), 97--178.

\bibitem{Yasaki}
D.~Yasaki, \emph{Perfect unary forms over real quadratic fields}, J. Th\'eor. Nombres Bordeaux
\textbf{25} (2013), 759--775.

\end{thebibliography}
\end{document}
